% Part: first-order-logic
% Chapter: syntax-and-semantics
% Section: assignments

\documentclass[../../../include/open-logic-section]{subfiles}

\begin{document}

\olfileid[pt]{fol}{syn}{ass}

\olsection{Atribuições de Variáveis}

\begin{explain}
Uma atribuição de variáveis~$s$ fornece um valor para \emph{toda}
variável---e há infinitas delas. Isso, é claro,
não é necessário. Exigimos que as atribuições de variáveis atribuam valores a
todas as !!{variable}s simplesmente porque isso torna as coisas muito mais fáceis.  O
valor de um termo~$t$, e se !!a{formula}~$!A$ é
satisfeita em !!a{structure} com respeito a~$s$, dependem apenas das
atribuições que~$s$ faz às !!{variable}s em~$t$ e às
!!{variable}s livres de~$!A$.  Esse é o conteúdo das próximas duas
proposições.  Para tornar precisa a ideia de ``depende de'', mostramos
que quaisquer duas atribuições de variáveis que concordam em todas as variáveis
em~$t$ dão o mesmo valor, e que $!A$ é satisfeita relativamente a uma
sse é satisfeita relativamente à outra se duas atribuições de variáveis
concordam em todas as variáveis livres de~$!A$.
\end{explain}

\begin{prop}\ollabel{prop:valindep}
Se as !!{variable}s em um termo~$t$ estão entre $x_1$, \dots,~$x_n$, e
$s_1(x_i) = s_2(x_i)$ para $i = 1$, \dots,~$n$, então $\Value{t}{M}[s_1]
= \Value{t}{M}[s_2]$.
\end{prop}

\begin{proof}
Por indução sobre a complexidade de~$t$. Para o caso base, $t$ pode ser
!!a{constant} ou uma das variáveis~$x_1$, \dots,~$x_n$.  Se $t
= c$, então $\Value{t}{M}[s_1] = \Assign{c}{M} = \Value{t}{M}[s_2]$. Se
$t = x_i$, $s_1(x_i) = s_2(x_i)$ pela hipótese da proposição,
e assim $\Value{t}{M}[s_1] = s_1(x_i) = s_2(x_i) = \Value{t}{M}[s_2]$.

Para o passo indutivo, suponha que $t = \Atom{f}{t_1, \dots, t_k}$
e que a afirmação vale para $t_1$, \dots, $t_k$. Então
\begin{align*}
  \Value{t}{M}[s_1] & = \Value{\Atom{f}{t_1, \dots, t_k}}{M}[s_1] \\
  & = \Assign{f}{M}(\Value{t_1}{M}[s_1], \dots, \Value{t_k}{M}[s_1]).
\intertext{Para $j = 1$, \dots,~$k$, as !!{variable}s de~$t_j$ estão
    entre $x_1$, \dots,~$x_n$. Pela hipótese de indução,
    $\Value{t_j}{M}[s_1] = \Value{t_j}{M}[s_2]$. Assim,}
\Value{t}{M}[s_1] & = \Value{\Atom{f}{t_1, \dots, t_k}}{M}[s_1] \\
  & = \Assign{f}{M}(\Value{t_1}{M}[s_1], \dots, \Value{t_k}{M}[s_1]) \\
  & = \Assign{f}{M}(\Value{t_1}{M}[s_2], \dots, \Value{t_k}{M}[s_2]) \\
  & = \Value{\Atom{f}{t_1, \dots, t_k}}{M}[s_2] = \Value{t}{M}[s_2].
\end{align*}
\end{proof}

\begin{prop}\ollabel{prop:satindep}
Se as !!{variable}s livres em $!A$ estão entre $x_1$, \dots,~$x_n$, e
$s_1(x_i) = s_2(x_i)$ para $i = 1$, \dots,~$n$, então $\Sat{M}{!A}[s_1]$
sse $\Sat{M}{!A}[s_2]$.
\end{prop}

\begin{proof}
Usamos indução sobre a complexidade de $!A$. Para o caso base, onde
$!A$ é atômica, $!A$ pode ser:
\iftag{prvTrue}{$\ltrue$,}{}
\iftag{prvFalse}{$\lfalse$,}{}
$\Atom{R}{t_1, \dots, t_k}$ para um predicado $R$ de $k$ lugares e termos
$t_1$, \dots,~$t_k$, ou $\eq[t_1][t_2]$ para termos $t_1$ e~$t_2$.
Nos dois últimos casos, demonstramos apenas a direção direta do
!!{biconditional}, já que a prova da recíproca é simétrica.

\begin{enumerate}
\tagitem{prvTrue}{%
  \indcase{!A}{\ltrue}{tanto $\Sat{M}{\indfrm}[s_1]$ quanto
    $\Sat{M}{\indfrm}[s_2]$.}}{}

\tagitem{prvFalse}{%
  \indcase{!A}{\lfalse}{tanto $\Sat/{M}{!A}[s_1]$ quanto
    $\Sat/{M}{!A}[s_2]$.}}{}

\item
  \indcase{!A}{\Atom{R}{t_1, \ldots, t_k}}{seja
    $\Sat{M}{\indfrm}[s_1]$. Então
    \[
    \langle \Value{t_1}{M}[s_1], \ldots, \Value{t_k}{M}[s_1] \rangle
    \in \Assign{R}{M}.
    \]
    Para $i = 1$, \dots,~$k$, $\Value{t_i}{M}[s_1] =
    \Value{t_i}{M}[s_2]$ por \olref{prop:valindep}. Então também temos
    $\langle \Value{t_i}{M}[s_2], \ldots, \Value{t_k}{M}[s_2] \rangle
    \in \Assign{R}{M}$, e portanto $\Sat{M}{\indfrm}[s_2]$.}
\item
  \indcase{!A}{\eq[t_1][t_2]}{suponha $\Sat{M}{\indfrm}[s_1]$.
    Então $\Value{t_1}{M}[s_1] = \Value{t_2}{M}[s_1]$. Assim,
    \begin{align*}
      \Value{t_1}{M}[s_2] & = \Value{t_1}{M}[s_1]
      & \text{(por \olref{prop:valindep})} \\
      & = \Value{t_2}{M}[s_1] 
      & \text{(já que $\Sat{M}{\eq[t_1][t_2]}[s_1]$)}\\
      &= \Value{t_2}{M}[s_2]
      & \text{(por \olref{prop:valindep}),}
    \end{align*}
    então $\Sat{M}{\eq[t_1][t_2]}[s_2]$.}
\end{enumerate}

Agora suponha $\Sat{M}{!B}[s_1]$ sse $\Sat{M}{!B}[s_2]$ para todas as
!!{formula}s $!B$ menos complexas que~$!A$. O passo de indução procede
por casos determinados pelo operador principal de~$!A$. Em cada caso,
demonstramos apenas a direção direta do !!{biconditional}; a
prova da direção reversa é simétrica. Em todos os casos exceto
os dos quantificadores, aplicamos a hipótese de indução às
sub-!!{formula}s~$!B$ de~$!A$.  As variáveis livres de~$!B$ estão entre
as de~$!A$. Assim, se $s_1$ e $s_2$ concordam nas variáveis livres
de~$!A$, elas também concordam nas de~$!B$, e a hipótese de indução
se aplica a~$!B$.

\begin{enumerate}
\tagitem{defNot}{}{%
  \iftag{probNot}{%
    \indcase!{!A}{\lnot !B}{}}{%
    \indcase{!A}{\lnot !B}{se $\Sat{M}{\indfrm}[s_1]$, então
      $\Sat/{M}{!B}[s_1]$, então pela hipótese de indução,
      $\Sat/{M}{!B}[s_2]$, portanto $\Sat{M}{\indfrm}[s_2]$.}}}

\tagitem{defAnd}{}{%
  \iftag{probAnd}{%
    \indcase!{!A}{!B \land !C}{}}{%
    \indcase{!A}{!B \land !C}{se $\Sat{M}{\indfrm}[s_1]$, então
      $\Sat{M}{!B}[s_1]$ e $\Sat{M}{!C}[s_1]$, então pela hipótese de
      indução, $\Sat{M}{!B}[s_2]$ e $\Sat{M}{!C}[s_2]$. Portanto,
      $\Sat{M}{\indfrm}[s_2]$.}}}

\tagitem{defOr}{}{%
  \iftag{probOr}{%
    \indcase!{!A}{!B \lor !C}{}}{%
    \indcase{!A}{!B \lor !C}{se $\Sat{M}{\indfrm}[s_1]$, então
      $\Sat{M}{!B}[s_1]$ ou $\Sat{M}{!C}[s_1]$. Pela hipótese de indução,
      $\Sat{M}{!B}[s_2]$ ou $\Sat{M}{!C}[s_2]$, então $\Sat{M}{\indfrm}[s_2]$.}}}

\tagitem{defIf}{}{%
  \iftag{probIf}{%
    \indcase!{!A}{!B \lif !C}{}}{%
    \indcase{!A}{!B \lif !C}{se $\Sat{M}{\indfrm}[s_1]$, então
    $\Sat/{M}{!B}[s_1]$ ou $\Sat{M}{!C}[s_1]$. Pela hipótese de indução,
    $\Sat/{M}{!B}[s_2]$ ou $\Sat{M}{!C}[s_2]$, então $\Sat{M}{\indfrm}[s_2]$.}}}

\tagitem{defIff}{}{%
  \iftag{probIff}{%
    \indcase!{!A}{!B \liff !C}{}}{%
    \indcase{!A}{!B \liff !C}{se $\Sat{M}{\indfrm}[s_1]$, então ou
      $\Sat{M}{!B}[s_1]$ e $\Sat{M}{!C}[s_1]$, ou $\Sat/{M}{!B}[s_1]$ e
      $\Sat/{M}{!C}[s_1]$. Pela hipótese de indução, ou
      $\Sat{M}{!B}[s_2]$ e $\Sat{M}{!C}[s_2]$ ou $\Sat/{M}{!B}[s_2]$
      e $\Sat/{M}{!C}[s_2]$. Em qualquer caso, $\Sat{M}{\indfrm}[s_2]$.}}}

\tagitem{defEx}{}{%
  \iftag{probEx}{%
    \indcase!{!A}{\lexists[x][!B]}{}}{%
    \indcase{!A}{\lexists[x][!B]}{se $\Sat{M}{\indfrm}[s_1]$, há
      um $m \in \Domain{M}$ tal que $\Sat{M}{!B}[\Subst{s_1}{m}{x}]$.
      Sejam $s_1' = \Subst{s_1}{m}{x}$ e $s_2' =
      \Subst{s_2}{m}{x}$. As variáveis livres de~$!B$ estão entre $x_1$,
      \dots, $x_n$, e $x$. $s_1'(x_i) = s_2'(x_i)$, já que $s_1'$ e
      $s_2'$ são $x$-variantes de $s_1$ e~$s_2$, respectivamente, e por
      hipótese $s_1(x_i) = s_2(x_i)$. $s_1'(x) = s_2'(x) = m$ pela
      forma como definimos $s_1'$ e~$s_2'$. Então a hipótese de
      indução se aplica a $!B$ e $s_1'$, $s_2'$, então
      $\Sat{M}{!B}[s_2']$. Portanto, já que $s_2' = \Subst{s_2}{m}{x}$,
      há um $m \in \Domain{M}$ tal que
      $\Sat{M}{!B}[\Subst{s_2}{m}{x}]$, e assim
      $\Sat{M}{\indfrm}[s_2]$.}}}

\tagitem{defAll}{}{%
  \iftag{probAll}{%
    \indcase!{!A}{\lforall[x][!B]}{}}{%
    \indcase{!A}{\lforall[x][!B]}{se $\Sat{M}{\indfrm}[s_1]$, então para
      todo $m \in \Domain{M}$, $\Sat{M}{!B}[\Subst{s_1}{m}{x}]$. Queremos
      mostrar que também, para todo $m \in \Domain{M}$,
      $\Sat{M}{!B}[\Subst{s_2}{m}{x}]$. Então seja $m \in \Domain{M}$
      arbitrário, e considere $s_1' = \Subst{s}{m}{x}$ e $s_2' =
      \Subst{s}{m}{x}$. Temos que $\Sat{M}{!B}[s_1']$. As variáveis
      livres de~$!B$ estão entre $x_1$, \dots, $x_n$, e $x$.
      $s_1'(x_i) = s_2'(x_i)$, já que $s_1'$ e $s_2'$ são
      $x$-variantes de $s_1$ e~$s_2$, respectivamente, e por hipótese
      $s_1(x_i) = s_2(x_i)$. $s_1'(x) = s_2'(x) = m$ pela forma como
      definimos $s_1'$ e~$s_2'$. Então a hipótese de indução
      se aplica a~$!B$ e~$s_1'$, $s_2'$, e temos
      $\Sat{M}{!B}[s_2']$. Isso se aplica a todo $m \in \Domain{M}$,
      isto é, $\Sat{M}{!B}[\Subst{s_2}{m}{x}]$ para todo~$m \in
      \Domain{M}$, então $\Sat{M}{\indfrm}[s_2]$.}}}
\end{enumerate}
Por indução, obtemos que $\Sat{M}{!A}[s_1]$ sse $\Sat{M}{!A}[s_2]$
sempre que as !!{variable}s livres em $!A$ estão entre $x_1$, \dots, $x_n$
e $s_1(x_i)=s_2(x_i)$ para $i = 1$, \dots,~$n$.
\end{proof}

\begin{probtag}{probNot,probOr,probAnd,probIf,probIff,probEx,probAll}
Complete a prova de \olref[fol][syn][ass]{prop:satindep}.
\end{probtag}

\begin{explain}
!!^{sentence}s não têm variáveis livres, então quaisquer duas atribuições de variáveis
atribuem as mesmas coisas a todas as (zero) variáveis livres de qualquer
sentença. A proposição que acabamos de provar significa, então, que se
!!a{sentence} é satisfeita ou não em uma estrutura relativamente a uma atribuição
de variáveis é completamente independente da atribuição. Registraremos
esse fato. Ele justifica a definição de satisfação de
!!a{sentence} em !!a{structure} (sem mencionar uma atribuição de
variáveis) que segue.
\end{explain}

\begin{cor}
\ollabel{cor:sat-sentence}
Se $!A$ é !!a{sentence} e $s$ uma atribuição de variáveis, então
$\Sat{M}{!A}[s]$ sse $\Sat{M}{!A}[s']$ para toda atribuição de
variáveis~$s'$.
\end{cor}

\begin{proof}
  Seja $s'$ uma atribuição de variáveis qualquer. Como $!A$ é !!a{sentence}, ela
  não tem variáveis livres, e assim toda atribuição de variáveis~$s'$
  trivialmente atribui as mesmas coisas a todas as variáveis livres de~$!A$ que
  faz~$s$. Então a condição de \olref{prop:satindep} é satisfeita,
  e temos $\Sat{M}{!A}[s]$ sse $\Sat{M}{!A}[s']$.
\end{proof}

\begin{defn}
\ollabel{defn:satisfaction}
Se $!A$ é !!a{sentence}, dizemos que !!a{structure}~$\Struct M$
\emph{satisfaz}~$!A$, $\Sat{M}{!A}$, sse $\Sat{M}{!A}[s]$ para todas as
atribuições de variáveis~$s$.
\end{defn}

Se $\Sat{M}{!A}$, também dizemos simplesmente que \emph{$!A$ é verdadeira
em~$\Struct{M}$.} A noção de satisfação se estende naturalmente
de !!{sentence}s individuais para conjuntos de !!{sentence}s.

\begin{defn}
\ollabel{defn:sat}
Se $\Gamma$ é um conjunto de !!{sentence}s~$\Gamma$, dizemos que
!!a{structure}~$\Struct M$ \emph{satisfaz}~$\Gamma$,
$\Sat{M}{\Gamma}$, sse $\Sat{M}{!A}$ para todo $!A \in \Gamma$.
\end{defn}

\begin{prop}\ollabel{prop:sentence-sat-true}
  Seja $\Struct{M}$ !!a{structure}, $!A$ !!a{sentence}, e $s$ uma
  atribuição de variáveis.  $\Sat{M}{!A}$ sse $\Sat{M}{!A}[s]$.
\end{prop}

\begin{proof}
Exercício.
\end{proof}

\begin{prob}
Prove \olref[fol][syn][ass]{prop:sentence-sat-true}
\end{prob}

\begin{prop}\ollabel{prop:sat-quant}
  Suponha que $!A(x)$ contém apenas $x$ livre, e $\Struct M$ é
  !!a{structure}. Então:
  \begin{tagenumerate}{prvEx,prvAll}
    \tagitem{prvEx}{$\Sat{M}{\lexists[x][!A(x)]}$ sse $\Sat{M}{!A(x)}[s]$
      para pelo menos uma atribuição de variáveis~$s$.}{}
    \tagitem{prvAll}{$\Sat{M}{\lforall[x][!A(x)]}$ sse $\Sat{M}{!A(x)}[s]$
  para todas as atribuições de variáveis~$s$.}{}
\end{tagenumerate}
\end{prop}

\begin{proof}
Exercício.
\end{proof}

\begin{prob}
Prove \olref[fol][syn][ass]{prop:sat-quant}.
\end{prob}

\begin{prob}
\DeclareRobustCommand{\VDash}{\mathrel{||}\joinrel\Relbar}  
Suponha que $\Lang L$ é uma linguagem sem !!{function}s. Dada uma
!!{structure}~$\Struct M$, $c$ !!a{constant} e $a \in \Domain M$,
defina $\Struct M[a/c]$ como a !!{structure} que é exatamente
como~$\Struct M$, exceto que $\Assign{c}{M[a/c]} = a$. Defina
$\Struct M \VDash !A$ para !!{sentence}s~$!A$ por:
\begin{enumerate}
\tagitem{prvFalse}{%
  \indcase{!A}{\lfalse}{não $\Struct{M} \VDash \indfrm$.}}{}

\tagitem{prvTrue}{%
  \indcase{!A}{\ltrue}{$\Struct{M} \VDash \indfrm$.}}{}

\item \indcase{!A}{\Atom{R}{d_1, \dots, d_n}}{$\Struct M \VDash \indfrm$
  sse $\langle \Assign{d_1}{M}, \dots, \Assign{d_n}{M} \rangle \in
  \Assign{R}{M}$.}

\item \indcase{!A}{\eq[d_1][d_2]}{$\Struct M \VDash \indfrm$ sse
  $\Assign{d_1}{M} = \Assign{d_2}{M}$.}

\tagitem{prvNot}{%
  \indcase{!A}{\lnot !B}{$\Struct{M} \VDash \indfrm$ sse
    não $\Struct{M} \VDash {!B}$.}}{}

\tagitem{prvAnd}{%
  \indcase{!A}{(!B \land !C)}{$\Struct{M} \VDash
    \indfrm$ sse $\Struct{M} \VDash!B$ e $\Struct{M} \VDash !C$.}}{}

\tagitem{prvOr}{%
  \indcase{!A}{(!B \lor !C)}{$\Struct{M} \VDash \indfrm$ sse
    $\Struct{M} \VDash !B$ ou $\Struct{M} \VDash !C$ (ou ambos).}}{}

\tagitem{prvIf}{%
  \indcase{!A}{(!B \lif !C)}{$\Struct{M} \VDash \indfrm$ sse
    não $\Struct {M} \VDash !B$ ou $\Struct M \VDash !C$ (ou ambos).}}{}

\tagitem{prvIff}{%
  \indcase{!A}{(!B \liff !C)}{$\Struct{M} \VDash {\indfrm}$ sse ou
    ambos $\Struct{M} \VDash {!B}$ e $\Struct{M} \VDash {!C}$, ou
    nem $\Struct{M} \VDash {!B}$ nem $\Struct{M} \VDash {!C}$.}}{}

\tagitem{prvAll}{%
  \indcase{!A}{\lforall[x][!B]}{$\Struct{M} \VDash {\indfrm}$ sse para
    todo $a \in \Domain{M}$, $\Struct{M[a/c]} \VDash \Subst{!B}{c}{x}$,
    se $c$ não ocorre em~$!B$.}}{}

\tagitem{prvEx}{%
  \indcase{!A}{\lexists[x][!B]}{$\Struct{M} \VDash
  {\indfrm}$ sse há um $a \in \Domain M$ tal que
  $\Struct{M[a/c]} \VDash \Subst{!B}{c}{x}$, se $c$ não ocorre
  em~$!B$.}}{}
\end{enumerate}
Sejam $x_1$, \dots, $x_n$ todas as !!{variable}s livres em~$!A$,
$c_1$, \dots, $c_n$ símbolos de constante não presentes em~$!A$,
$a_1$, \dots, $a_n \in \Domain M$, e $s(x_i) = a_i$.

Mostre que $\Sat{M}{!A}[s]$ sse $\Struct M[a_1/c_1,\dots,a_n/c_n]
\VDash \Subst{\Subst{!A}{c_1}{x_1}\dots}{c_n}{x_n}$.

(Este problema mostra que é possível dar uma semântica para a
lógica de primeira ordem que dispensa atribuições de variáveis.)
\end{prob}

\begin{prob}
Suponha que $f$ é um símbolo de função não presente em~$!A(x,y)$. Mostre que
há !!a{structure}~$\Struct{M}$ tal que
$\Sat{M}{\lforall[x][\lexists[y][!A(x,y)]]}$ sse há uma~$\Struct
M'$ tal que $\Sat{M'}{\lforall[x][!A(x,f(x))]}$.

(Este problema é um caso especial do que é conhecido como Teorema de Skolem;
$\lforall[x][!A(x,f(x))]$ é chamado de \emph{forma normal de Skolem} de
$\lforall[x][\lexists[y][!A(x,y)]]$.)
\end{prob}

\end{document}
